\section{Ordinary group-word recovery}\label{sec:group-words}

The scalar cover bounds also apply to a nonlinear class with arbitrarily
large finite alphabets. An ordinary group word uses variables, their
inverses, multiplication, and fixed constants. The expression at a
coordinate is fixed on the whole code; arbitrary symbol maps and
endomorphisms are not included here. The latter require the vector-layer
benchmark of~\cite[Theorem B.5]{BrendeDirected}.

\begin{lemma}[The scalar exponent check]\label{lem:word-check}
\notformalized
Let $G$ be a finite $p$-group and $\mathcal D\subseteq G^N$ a set
satisfying fixed ordinary-word equations $x_v=w_v(x)$.
Let $E_{vu}\in\F_p$ be the signed exponent sum of variable $x_u$
in $w_v$, and put $k=\dim_{\F_p}\ker(I-E)$. Then
\[
 |\mathcal D|\le |G|^k.
\]
For any law on a nonempty $\mathcal D$, its entropy is at most
$k\log_2|G|$.
\end{lemma}
\begin{informalproof}
Use the terminating lower exponent-$p$ central series
$P_1=G$, $P_{j+1}=P_j^p[P_j,G]$. Its termination and the central-layer
argument are proved in~\cite[Theorem B.5]{BrendeDirected}.
Write $t_j=\dim_{\F_p}(P_j/P_{j+1})$, so
$\sum_jt_j=\log_p|G|$. Fix an attained solution modulo $P_j$ and
choose representatives modulo $P_{j+1}$. Each further lift has the
form $r_va_v$, where $a_v$ is in the central elementary abelian layer.
Centrality makes the change in an ordinary word equal to
$\sum_u E_{vu}a_u$: products add, inversion negates, and constants
give zero. Thus the lift equations have linear part
$(I-E)\otimes I_{t_j}$. A nonempty fibre is contained in an affine
solution space of size $p^{kt_j}$. Multiplying the fibre bounds gives
$|\mathcal D|\le p^{k\sum_jt_j}=|G|^k$. Empty sets and trivial groups
satisfy the cardinality bound directly. Entropy is at most the
logarithm of the support size.
\end{informalproof}

This is the ordinary-word specialization of the classical central
lifting method discussed in~\cite{BrendeDirected}, not a new general
group-equation theorem. In particular, a zero exponent sum is a
statement about this word language, not a rule for erasing arbitrary
nonlinear operations from a recovery circuit.

\begin{corollary}[Uniform bounds for growing group alphabets]
\label{cor:group-cover}
\notformalized
Fix the template, prime, embedding, chain, and positive integer type
widths of \cref{thm:fixed-width}. At level $i$, use any nontrivial
finite $p$-group $G_i$ as the alphabet in each of the
$N_i=Bq_i$ symbol positions. Let a nonempty code
$\mathcal D_i\subseteq G_i^{N_i}$ have ordinary-word recovery from
the permitted parents. Suppose its exponent matrix modulo $p$ is
homogeneous on the external sheets, with the block shapes of the
fixed template. For every $\varepsilon>0$, eventually
\[
 \frac{\log_{|G_i|}|\mathcal D_i|}{N_i}<\frac12+\varepsilon,
\]
with one threshold independent of $G_i$, its order, the words, and
their constants.

On the independent-chord family of \cref{thm:weighted-zero}, use one
$G_i$-symbol per vertex. If the exponent on each directed edge is
homogeneous and nonzero modulo $p$, the same normalized rate tends
to zero, uniformly in these choices. The corresponding entropy
bounds hold for every law on each code.
\end{corollary}
\begin{informalproof}
The exponent check $I-E_i$ is a characteristic-$p$ homogeneous
linear check with the stated fixed type widths. Apply
\cref{lem:word-check} to bound the normalized logarithmic size by
$\dim_{\F_p}\ker(I-E_i)/N_i$, then apply
\cref{thm:fixed-width}. For the second conclusion, the diagonal is
$1$ and the directed edge coefficients are the negatives of the
nonzero exponent sums. They meet \cref{thm:weighted-zero} exactly.
Its uniform threshold gives the zero-rate conclusion. The entropy
assertions follow from the final part of \cref{lem:word-check}.
\end{informalproof}

The words themselves may vary between sheets: only their exponent
matrices must be homogeneous. Reordering occurrences, changing
constants, and inserting commutators can preserve this matrix.
The number of group-symbol slots per type remains fixed, even though
the group alphabet may grow. Arbitrary endomorphic expressions have
matrix-valued central derivatives; the scalar argument above does
not cover them. Nor does it supply a group-word representation or
a charged native implementation for an arbitrary tape fibre.

This extends the cover obstruction beyond linear codes while retaining
an explicit algebraic restriction on recovery. An above-half candidate
in this cover model must violate the word or exponent hypotheses, or
the cover hypotheses themselves. The rate-zero conclusion additionally
requires every prescribed edge exponent to be nonzero; the single-edge
example of \cref{sec:weighted} shows why that condition matters.
